% This thesis studies the training dynamics of overparameterized neural networks
% through spectral bias, the tendency of gradient-based training to learn
% low-frequency components of a target function before high-frequency ones. The
% setting is regression on the unit circle with a one-hidden-layer ReLU network and
% targets given by finite Fourier mixtures, chosen because it is simple enough for
% the dynamics to be written down explicitly yet rich enough to exhibit spectral
% bias.
%
% In the infinite-width limit the residual dynamics are governed by the neural
% tangent kernel. Under the uniform measure on the circle this operator is a
% convolution, so it is diagonalised by the Fourier basis, each frequency decays at
% a rate set by its own eigenvalue, and low frequencies are learned faster because
% their eigenvalues are larger. This gives an explicit reference model for spectral
% bias against which finite effects can be measured.
%
% We then ask how much of this picture survives away from the idealised limit,
% working throughout on a fixed low-frequency Fourier subspace. Replacing the
% uniform measure by finitely many sampled points, we prove non-asymptotic bounds
% showing that the sampled operator stays close to the continuum Fourier prediction,
% with error of order \(O(n^{-1/2})\) in the sample size. Freezing the tangent
% kernel at finite width, we prove an analogous concentration around the
% infinite-width operator at rate \(O(m^{-1/2})\) in the width. In both cases the
% Fourier description degrades gracefully rather than breaking, and the perturbation
% shrinks as the sample size or width grows.
%
% The fixed-kernel account stops once the kernel is allowed to evolve during
% training. Here our analysis is empirical: at small width the evolving kernel
% reaches a lower loss than the frozen one, and it does so not by aligning its
% eigenspaces more closely with the Fourier basis, which can in fact degrade, but by
% increasing the spectral strength of the lowest-frequency blocks. Training therefore
% reshapes the kernel in a direction that reinforces the low-frequency bias rather
% than merely approximating the fixed-kernel dynamics, an effect that diminishes as
% width grows and the kernel moves less. We give evidence that this reweighting is
% tied to the network's parameterisation, since the constant and first-frequency
% modes are the ones the input embedding and bias represent most directly. A formal
% theory of this evolving-kernel regime remains the main problem left open.

In this thesis, we study spectral bias, the tendency of gradient-based training to
learn the low-frequency part of a target before its high-frequency part. We work
in a setting simple enough to analyse explicitly: regression on the unit circle
with a shallow ReLU network. In the infinite-width limit, the residual dynamics
are governed by the Neural Tangent Kernel. Under the uniform measure on the
circle this kernel depends only on the angle between two points, so the
associated operator is a convolution and the Fourier modes are its
eigenfunctions, each decaying at a rate set by its eigenvalue, and the lower the
frequency, the larger the eigenvalue, so low frequencies are learned first.

Away from this idealised limit the picture degrades only gradually. On a fixed
low-frequency subspace, both finite sampling and frozen finite width keep the
operator close to the continuum Fourier prediction, with error of order
\(O(n^{-1/2})\) in the sample size \(n\) and \(O(m^{-1/2})\) in the width \(m\).
The description breaks only once the kernel is allowed to evolve during training.
At small width the evolving kernel reaches a lower loss by strengthening its
lowest-frequency components, even as its alignment with the Fourier basis fails
to improve. This reinforces the low-frequency bias rather than approximating the
fixed-kernel dynamics. A formal theory of this evolving-kernel regime remains the
main open problem.
